\documentclass[11pt]{article}

\usepackage[margin=1in]{geometry}
\usepackage{amsmath,amssymb,amsthm,mathtools,bm}
\usepackage{mathrsfs}
\usepackage{bbm}
\usepackage{microtype}
\usepackage{xcolor}

\title{Global Convergence for Low-Rank Mean-Field Neural Networks: Proof Outline}
\author{ }
\date{}

% math macros
\newcommand{\E}{\mathbb{E}}
\newcommand{\Pp}{\mathbb{P}}
\newcommand{\R}{\mathbb{R}}
\newcommand{\1}{\mathbbm{1}}
\newcommand{\ip}[2]{\left\langle #1,\,#2 \right\rangle}
\newcommand{\norm}[1]{\left\lVert #1 \right\rVert}
\newcommand{\llbracket}[1]{[\![ #1 ]\!]}

\newtheorem{theorem}{Theorem}
\newtheorem{lemma}{Lemma}
\newtheorem{definition}{Definition}
\newtheorem{assumption}{Assumption}
\newtheorem{remark}{Remark}
\newtheorem{proposition}{Proposition}

\begin{document}

\maketitle

\section{Setup: Low-Rank Mean-Field Equations}

We consider a three-layer low-rank random feature network with the following structure:
\begin{itemize}
\item \textbf{First layer}: Random features $\varphi_1(\ip{f_1(C_1)}{X})$ where $f_1(C_1)$ are frozen random vectors and $C_1$ indexes the first-layer channels.
\item \textbf{Second layer}: Low-rank mixing matrix $L$ of rank $r$, producing outputs
\[
H_2(t,c_2;X,W)=\sum_{k=1}^r L_{c_2,k}\,m_k(t;X,W),
\]
where $m_k(t;X,W)=\E_{C_1}[w_1(t,C_1,k)\,\varphi_1(\ip{f_1(C_1)}{X})]$ are the $r$ channel-wise partial functions.
\item \textbf{Third layer}: Trainable weights $w_2(t,C_2)$ producing the final output.
\end{itemize}

The mean-field ODE system consists of $r+1$ equations:
\begin{align}
\partial_t w_1(t,c_1,k) &= -\xi_1(t)\,\E_{Z}\!\left[d_L\big(Z;W(t)\big)\,\varphi_1\!\left(\ip{f_1(c_1)}{X}\right)\,B_k(t;X)\right], \label{eq:ode-w1}\\
\partial_t w_2(t,c_2) &= -\xi_2(t)\,\E_{Z}\!\left[d_L\big(Z;W(t)\big)\,\varphi_2\!\Big(H_2(t,c_2;X,W(t))\Big)\right], \label{eq:ode-w2}
\end{align}
where $d_L(Z;W)=\partial_2 \mathcal{L}(Y,\hat y(X;W))$ is the loss derivative, and
\[
B_k(t;X)=\E_{C_2}\!\left[L_{C_2,k}\,\varphi_2'\!\Big(H_2(t,C_2;X,W(t))\Big)\,w_2(t,C_2)\right]
\]
is the channel-$k$ backpropagated signal.

The model output is:
\[
\hat y(X;W(t)) = \E_{C_2}\big[w_2(t,C_2)\,\varphi_2\big(H_2(t,C_2;X,W(t))\big)\big].
\]

\section{Assumptions for Global Convergence}

\begin{assumption}[Initialization]\label{assump:init-lowrank}
The initial functions $w_1^0:\Omega_1\times\{1,\dots,r\}\to\R$ and $w_2^0:\Omega_2\to\R$ satisfy:
\begin{align*}
\sup_{m\ge 1}\frac{1}{\sqrt{m}}\max_{1\le k\le r}\E_{C_1}\!\left[|w_1^0(C_1,k)|^m\right]^{1/m} &\le K,\\
\sup_{m\ge 1}\frac{1}{\sqrt{m}}\E_{C_2}\!\left[|w_2^0(C_2)|^m\right]^{1/m} &\le K.
\end{align*}
\end{assumption}

\begin{assumption}[Diversity]\label{assump:diversity-lowrank}
The support of $\rho^1$ (the measure on $\Omega_1$) is $\R^d$ (or dense in $\R^d$). This ensures that the random features $\varphi_1(\ip{f_1(C_1)}{\cdot})$ have dense span in $L^2(\mathcal{P}_X)$.
\end{assumption}

\begin{assumption}[Regularity]\label{assump:regularity-lowrank}
There exists $K\ge 1$ such that:
\begin{itemize}
\item $\xi_1,\xi_2$ are bounded on compact intervals and piecewise continuous.
\item $\varphi_1,\varphi_2,\varphi_2'$ are $K$-bounded and $K$-Lipschitz; moreover $\varphi_2'$ is non-vanishing (strictly non-zero everywhere).
\item $\partial_2\mathcal{L}(y,\cdot)$ is $K$-bounded and $K$-Lipschitz for all $y$ in the support of $\mathcal{P}$.
\item $X$ is $K$-sub-Gaussian and $\norm{f_1(c_1)}\le K$ for all $c_1$.
\item The mixing kernel satisfies $\sup_{c_2,k}|L_{c_2,k}|\le K$ and $k\mapsto L_{c_2,k}$ is measurable for each $c_2$.
\item $|X|\le K$ with probability 1.
\end{itemize}
\end{assumption}

\begin{assumption}[Convergence]\label{assump:convergence-lowrank}
There exist functions $\bar{w}_1:\Omega_1\times\{1,\dots,r\}\to\R$ and $\bar{w}_2:\Omega_2\to\R$ such that as $t\to\infty$, there exists a coupling $\pi_t$ of $\rho^1\times\rho^2$ and itself such that:
\begin{align}
\int (1+|\bar{w}_2(c_2)|)|\bar{w}_2(c_2)|\max_{1\le k\le r}|\bar{w}_1(c_1,k)|\,|w_1^*(t,c_1',k)-\bar{w}_1(c_1,k)|\,d\pi_t(c_1,c_2,c_1',c_2') &\to 0, \label{eq:conv-w1}\\
\int (1+|\bar{w}_2(c_2)|)|\bar{w}_2(c_2)|\,|w_2^*(t,c_2')-\bar{w}_2(c_2)|\,d\pi_t(c_1,c_2,c_1',c_2') &\to 0, \label{eq:conv-w2}
\end{align}
where $W^*(t)=(w_1^*(t,\cdot,\cdot),w_2^*(t,\cdot))$ is the solution to the mean-field ODEs \eqref{eq:ode-w1}--\eqref{eq:ode-w2}.
\end{assumption}

\begin{remark}
The convergence assumption \eqref{eq:conv-w1}--\eqref{eq:conv-w2} is adapted from the full-width case. The key difference is that we have $r$ channels for $w_1$, so we take the maximum over $k\in\{1,\dots,r\}$ in \eqref{eq:conv-w1}. The low-rank structure enters through the factor $\max_k|\bar{w}_1(c_1,k)|$ which is controlled by the $\psi_2$ bounds.
\end{remark}

\section{Key Technical Lemma: Bi-Lipschitz Property of $H_2$}

\begin{lemma}[Bi-Lipschitz property of $H_2$ with low-rank structure]\label{lem:bilipschitz-H2}
For any $W'=(w_1',w_2')$ and $W''=(w_1'',w_2'')$ satisfying the regularity assumptions, we have:
\[
\big|H_2(t,c_2;X,W')-H_2(t,c_2;X,W'')\big| \le K\,\|L\|_{\infty,1}\,\max_{1\le k\le r}\E_{C_1}\big[|w_1'(t,C_1,k)-w_1''(t,C_1,k)|\big],
\]
where $\|L\|_{\infty,1}\equiv \sup_{c_2}\sum_{k=1}^r |L_{c_2,k}|\le rK$ under the entrywise bound.

Moreover, if $\varphi_2$ is ReLU or a Lipschitz activation, then:
\[
\big|\varphi_2(H_2(t,c_2;X,W'))-\varphi_2(H_2(t,c_2;X,W''))\big| \le K\,\|L\|_{\infty,1}\,\max_{1\le k\le r}\E_{C_1}\big[|w_1'(t,C_1,k)-w_1''(t,C_1,k)|\big].
\]
\end{lemma}

\begin{proof}[Proof sketch]
By definition:
\begin{align*}
H_2(t,c_2;X,W')-H_2(t,c_2;X,W'') &= \sum_{k=1}^r L_{c_2,k}\big(m_k(t;X,W')-m_k(t;X,W'')\big)\\
&= \sum_{k=1}^r L_{c_2,k}\E_{C_1}\big[(w_1'(t,C_1,k)-w_1''(t,C_1,k))\,\varphi_1(\ip{f_1(C_1)}{X})\big].
\end{align*}
Taking absolute values and using boundedness of $\varphi_1$:
\begin{align*}
\big|H_2(t,c_2;X,W')-H_2(t,c_2;X,W'')\big| &\le \sum_{k=1}^r |L_{c_2,k}|\,\E_{C_1}\big[|w_1'(t,C_1,k)-w_1''(t,C_1,k)|\big]\,|\varphi_1(\ip{f_1(C_1)}{X})|\\
&\le K\,\sum_{k=1}^r |L_{c_2,k}|\,\max_{1\le k\le r}\E_{C_1}\big[|w_1'(t,C_1,k)-w_1''(t,C_1,k)|\big]\\
&\le K\,\|L\|_{\infty,1}\,\max_{1\le k\le r}\E_{C_1}\big[|w_1'(t,C_1,k)-w_1''(t,C_1,k)|\big].
\end{align*}
The second part follows from the Lipschitz property of $\varphi_2$.
\end{proof}

\section{Proof Outline: Global Convergence}

\begin{theorem}[Global Convergence for Low-Rank Networks]\label{thm:global-convergence-lowrank}
Under Assumptions~\ref{assump:init-lowrank}--\ref{assump:convergence-lowrank}, if the loss $\mathcal{L}$ is convex in its second variable, then the limit point $(\bar{w}_1,\bar{w}_2)$ is a global minimizer of the population loss:
\[
\mathscr{L}(\bar{w}_1,\bar{w}_2) = \E_Z\big[\mathcal{L}\big(Y,\hat y(X;\bar{w}_1,\bar{w}_2)\big)\big].
\]
\end{theorem}

\subsection{Detailed Proof (Adapted from Full-Width Case)}

We now provide the detailed proof, adapting the argument from the full-width three-layer case to our low-rank setting. The key simplification is that we do not train $\rho^1$ (the first layer uses frozen random features), so we avoid the algebraic topology arguments needed to show full support of $\text{Law}(w_1^*(t,C_1))$.

\textbf{Step 1: Convergence of the backpropagated signal}

By the convergence assumption (Assumption~\ref{assump:convergence-lowrank}), we have that for any $\epsilon>0$, there exists $T(\epsilon)$ such that for all $t\ge T(\epsilon)$ and almost surely:
\[
\epsilon \ge \left|\E_{Z}\!\left[d_L\big(Z;W^*(t)\big)\,\varphi_1\!\left(\ip{f_1(C_1)}{X}\right)\,B_k^*(t;X)\right]\right| = \left|\left\langle \E_{Z}\!\left[d_L\big(Z;W^*(t)\big)\,B_k^*(t;X)\middle|X=x\right],\varphi_1(\ip{f_1(C_1)}{x})\right\rangle_{L^2(\mathcal{P}_X)}\right|,
\]
where $B_k^*(t;X)=\E_{C_2}[L_{C_2,k}\,\varphi_2'(H_2^*(t,C_2;X,W^*(t)))\,w_2^*(t,C_2)]$.

Since $f_1(C_1)$ are frozen random features with full support in $\R^d$ (by Assumption~\ref{assump:diversity-lowrank}), and we have $r$ channels, the set $\{\varphi_1(\ip{f_1(c_1)}{\cdot}): c_1\in\Omega_1, k\in\{1,\dots,r\}\}$ has dense span in $L^2(\mathcal{P}_X)$. We thus obtain that for $u$ in a dense subset of $\R^d$:
\[
\text{ess-sup}\left|\left\langle \E_{Z}\!\left[d_L\big(Z;W^*(t)\big)\,B_k^*(t;X)\middle|X=x\right],\varphi_1(\ip{u}{x})\right\rangle_{L^2(\mathcal{P}_X)}\right| \le \epsilon.
\]

By continuity of $u\mapsto\varphi_1(\ip{u}{\cdot})$ in $L^2(\mathcal{P}_X)$, we extend the above to all $u\in\R^d$.

\textbf{Step 2: Coupling argument with low-rank structure}

Recall the couplings $\pi_t$ in Assumption~\ref{assump:convergence-lowrank}. Since $\varphi_1$ is bounded, we have:
\begin{align*}
&\E_{(C_1,C_2,C_1',C_2')\sim\pi_t}\left[\left|\left\langle \E_{Z}\!\left[d_L\big(Z;W^*(t)\big)\,B_k^*(t;X)-d_L\big(Z;\bar{W}\big)\,\bar{B}_k(X)\middle|X=x\right],\varphi_1(\ip{u}{x})\right\rangle_{L^2(\mathcal{P}_X)}\right|\right]\\
&\quad \le K\,\E_{\pi_t}\left[\left|d_L\big(Z;W^*(t)\big)\,B_k^*(t;X)-d_L\big(Z;\bar{W}\big)\,\bar{B}_k(X)\right|\right],
\end{align*}
where $\bar{W}=(\bar{w}_1,\bar{w}_2)$ and $\bar{B}_k(X)=\E_{C_2}[L_{C_2,k}\,\varphi_2'(\bar{H}_2(C_2;X,\bar{W}))\,\bar{w}_2(C_2)]$ with $\bar{H}_2(c_2;X,\bar{W})=\sum_{k=1}^r L_{c_2,k}\,\bar{m}_k(X,\bar{W})$ and $\bar{m}_k(X,\bar{W})=\E_{C_1}[\bar{w}_1(C_1,k)\,\varphi_1(\ip{f_1(C_1)}{X})]$.

\textbf{Step 3: Key bound using bi-Lipschitz property (adapted from full-width case)}

Using the regularity assumption and Lemma~\ref{lem:bilipschitz-H2}, similar to the calculation in the full-width case, we bound:
\begin{align*}
&\E_{\pi_t}\left[\left|d_L\big(Z;W^*(t)\big)\,B_k^*(t;X)-d_L\big(Z;\bar{W}\big)\,\bar{B}_k(X)\right|\right]\\
&\quad \le K\,\E_{\pi_t}\Big[(1+|\bar{w}_2(C_2)|)\Big(|w_2^*(t,C_2')-\bar{w}_2(C_2)|\\
&\qquad + |\bar{w}_2(C_2)|\,\max_{1\le k\le r}|w_1^*(t,C_1',k)-\bar{w}_1(C_1,k)|\Big)\Big],
\end{align*}
where the last step uses:
\begin{itemize}
\item The Lipschitz property of $\partial_2\mathcal{L}$ and $\varphi_2'$
\item The bi-Lipschitz bound from Lemma~\ref{lem:bilipschitz-H2} for $H_2$ differences: $|H_2^*(t)-H_2(\bar{W})|\le K\|L\|_{\infty,1}\max_k|w_1^*(t)-\bar{w}_1|$
\item The structure $B_k = \E_{C_2}[L_{C_2,k}\,\varphi_2'(H_2)\,w_2]$, which gives:
\begin{align*}
|B_k^*(t)-B_k(\bar{W})| &\le K\E_{C_2}\big[|L_{C_2,k}|\big(|\varphi_2'(H_2^*(t))-\varphi_2'(\bar{H}_2)|\,|w_2^*(t)| + |\varphi_2'(\bar{H}_2)|\,|w_2^*(t)-\bar{w}_2|\big)\big]\\
&\le K(1+|\bar{w}_2|)\Big(|w_2^*(t)-\bar{w}_2| + |\bar{w}_2|\,\max_k|w_1^*(t)-\bar{w}_1|\Big)
\end{align*}
\end{itemize}

By Assumption~\ref{assump:convergence-lowrank}, the right-hand side converges to $0$ as $t\to\infty$.

\textbf{Step 4: Orthogonality and optimality}

We thus obtain that for all $u\in\R^d$:
\[
\E_{C_2}\left[\left|\left\langle \E_{Z}\!\left[d_L\big(Z;\bar{W}\big)\,\bar{B}_k(X)\middle|X=x\right],\varphi_1(\ip{u}{x})\right\rangle_{L^2(\mathcal{P}_X)}\right|\right] = 0,
\]
which yields that for all $u\in\R^d$ and almost surely:
\[
\left|\left\langle \E_{Z}\!\left[d_L\big(Z;\bar{W}\big)\,\bar{B}_k(X)\middle|X=x\right],\varphi_1(\ip{u}{x})\right\rangle_{L^2(\mathcal{P}_X)}\right| = 0.
\]

Since $\{\varphi_1(\ip{u}{\cdot}): u\in\R^d\}$ has dense span in $L^2(\mathcal{P}_X)$ (by Assumption~\ref{assump:diversity-lowrank}), we have:
\[
\E_{Z}\!\left[d_L\big(Z;\bar{W}\big)\,\bar{B}_k(X)\middle|X=x\right] = 0
\]
for $\mathcal{P}_X$-almost every $x$ and almost every $c_2$.

Expanding $\bar{B}_k(X)$:
\[
\E_{Z}\!\left[\partial_2\mathcal{L}\big(Y,\hat y(X;\bar{W})\big)\middle|X=x\right]\,\E_{C_2}\!\left[L_{C_2,k}\,\varphi_2'(\bar{H}_2(C_2;X,\bar{W}))\,\bar{w}_2(C_2)\right] = 0.
\]

\textbf{Step 5: Non-degeneracy}

We note that our assumptions guarantee that $\Pp(\bar{w}_2(C_2)\ne 0)$ is positive. This is a crucial step that requires careful justification. In the full-width case, there are two cases:
\begin{itemize}
\item Case A: If $\int \mathbb{I}(u_2\ne 0)\rho^2(du_2)>0$ and $\xi_2(\cdot)=0$ (second layer not trained), then it is obvious that $\Pp(\bar{w}_2(C_2)\ne 0)>0$ since the initialization has non-zero mass.

\item Case B: If $\xi_2(\cdot)\ne 0$ (second layer is trained) and $\mathscr{L}(w_1^0,w_2^0) < \E_Z[\mathcal{L}(Y,\varphi_2(0))]$, then by the gradient flow property (loss decreases along trajectories):
\[
\mathscr{L}(w_1^*(t),w_2^*(t)) \le \mathscr{L}(w_1^*(t'),w_2^*(t')) \quad \text{for } t\ge t'.
\]
In particular, setting $t'=0$ and taking $t\to\infty$:
\[
\mathscr{L}(\bar{w}_1,\bar{w}_2) \le \mathscr{L}(w_1^0,w_2^0) < \E_Z[\mathcal{L}(Y,\varphi_2(0))].
\]
If $\Pp(\bar{w}_2(C_2)=0)=1$, then $\hat y(X;\bar{w}_1,\bar{w}_2)=\E_{C_2}[\bar{w}_2(C_2)\,\varphi_2(\bar{H}_2)]=0$ (since $\bar{w}_2=0$), so $\mathscr{L}(\bar{w}_1,\bar{w}_2)=\E_Z[\mathcal{L}(Y,\varphi_2(0))]$, a contradiction.
\end{itemize}

\begin{remark}[On the assumption $\mathscr{L}(w_1^0,w_2^0) < \E_Z[\mathcal{L}(Y,\varphi_2(0))]$]
This assumption may seem restrictive, but it is actually quite natural:
\begin{itemize}
\item It requires that the initial loss is better than the trivial predictor $\hat y=0$ (which corresponds to $w_2=0$).
\item For most reasonable initializations (e.g., small random weights), this condition is satisfied with high probability.
\item If this condition fails, then the network starts worse than the trivial solution, and convergence to a global minimizer (which must be at least as good as the trivial solution) is still possible, but the non-degeneracy argument needs to be modified.
\item In practice, this assumption can be verified by checking the initial loss value.
\end{itemize}
\end{remark}

Since $\varphi_2'$ is strictly non-zero and the mixing kernel $L$ has non-zero entries, we have:
\[
\E_{Z}\!\left[\partial_2\mathcal{L}\big(Y,\hat y(X;\bar{W})\big)\middle|X=x\right] = 0
\]
for $\mathcal{P}_X$-almost every $x$.

\textbf{Step 6: Global optimality}

Since $\mathcal{L}$ is convex in its second variable, for any measurable function $\tilde{y}(x)$:
\[
\mathcal{L}\big(y,\tilde{y}(x)\big)-\mathcal{L}\big(y,\hat y(x;\bar{W})\big) \ge \partial_2\mathcal{L}\big(y,\hat y(x;\bar{W})\big)\big(\tilde{y}(x)-\hat y(x;\bar{W})\big).
\]

Taking expectation:
\[
\E_Z\big[\mathcal{L}\big(Y,\tilde{y}(X)\big)\big] \ge \mathscr{L}(\bar{w}_1,\bar{w}_2),
\]
i.e., $(\bar{w}_1,\bar{w}_2)$ is a global minimizer.

\textbf{Step 7: Connection to the limit (exact adaptation)}

Finally, to connect $\mathscr{L}(W^*(t))$ with $\mathscr{L}(\bar{w}_1,\bar{w}_2)$ in the limit $t\to\infty$, we have (exactly as in the full-width case):
\begin{align*}
\big|\mathscr{L}(W^*(t))-\mathscr{L}(\bar{w}_1,\bar{w}_2)\big| &= \left|\E_Z\big[\mathcal{L}\big(Y,\hat y^*(t,X)\big)-\mathcal{L}\big(Y,\hat y(X;\bar{W})\big)\big]\right|\\
&\le K\,\E_Z\big[\big|\hat y^*(t,X)-\hat y(X;\bar{W})\big|\big]\\
&\le K\,\E_{\pi_t}\Big[|w_2^*(t,C_2')-\bar{w}_2(C_2)|\\
&\quad + |\bar{w}_2(C_2)|\,\max_{1\le k\le r}|w_1^*(t,C_1',k)-\bar{w}_1(C_1,k)|\Big],
\end{align*}
where the last step uses:
\begin{itemize}
\item $\hat y(X;W) = \E_{C_2}[w_2(C_2)\,\varphi_2(H_2(C_2;X,W))]$
\item The Lipschitz property of $\varphi_2$ and the bi-Lipschitz bound for $H_2$ from Lemma~\ref{lem:bilipschitz-H2}
\item The structure: $|\hat y^*(t)-\hat y(\bar{W})|\le K\E_{C_2}[|w_2^*(t)-\bar{w}_2| + |\bar{w}_2|\,|H_2^*(t)-\bar{H}_2|]$
\end{itemize}

By Assumption~\ref{assump:convergence-lowrank}, the right-hand side converges to $0$ as $t\to\infty$. This completes the proof.

\section{Key Differences from Full-Width Case}

\begin{enumerate}
\item \textbf{Simplification: No algebraic topology needed}: Since we do not train $\rho^1$ (the first layer uses frozen random features), we avoid the complex arguments needed to show that $\text{Law}(w_1^*(t,C_1))$ has full support. Instead, we rely on the fact that $f_1(C_1)$ are frozen random features with full support, which directly gives the dense span property.

\item \textbf{Low-rank structure in $H_2$}: The pre-activation $H_2$ is a sum over $r$ channels, introducing a factor $\|L\|_{\infty,1}\le rK$ in the bi-Lipschitz bound (Lemma~\ref{lem:bilipschitz-H2}).

\item \textbf{Convergence assumption}: Instead of bounding differences for a single $w_1$, we take the maximum over $r$ channels: $\max_{1\le k\le r}|w_1^*(t,c_1',k)-\bar{w}_1(c_1,k)|$. This is the key adaptation: the convergence assumption becomes:
\[
\int (1+|\bar{w}_2|)|\bar{w}_2|\,\max_{1\le k\le r}|\bar{w}_1(c_1,k)|\,|w_1^*(t,c_1',k)-\bar{w}_1(c_1,k)|\,d\pi_t \to 0.
\]

\item \textbf{Backpropagated signal}: The signal $B_k$ involves the mixing kernel $L$, but the structure remains similar: $B_k = \E_{C_2}[L_{C_2,k}\,\varphi_2'(H_2)\,w_2]$. The bound becomes:
\[
|B_k^*(t)-B_k(\bar{W})| \le K(1+|\bar{w}_2|)\Big(|w_2^*(t)-\bar{w}_2| + |\bar{w}_2|\,\max_k|w_1^*(t)-\bar{w}_1|\Big).
\]

\item \textbf{Universal approximation}: The dense span property still holds because we have $r$ channels with frozen random features $f_1(C_1)$ having full support. The low-rank structure does not reduce the expressivity below what is needed for global convergence, as long as $r$ is sufficiently large.
\end{enumerate}

\section{Discussion on Assumptions}

\subsection{The Non-Degeneracy Assumption}

The assumption that $\mathscr{L}(w_1^0,w_2^0) < \E_Z[\mathcal{L}(Y,\varphi_2(0))]$ (initial loss better than trivial predictor) is used to ensure non-degeneracy of the limit point. While this may seem restrictive, it is actually quite natural:

\begin{itemize}
\item \textbf{Practical relevance}: For most reasonable initializations (e.g., small random weights, Xavier initialization), this condition is satisfied with high probability. The initial network typically performs better than the trivial solution $\hat y=0$.

\item \textbf{Verifiability}: This assumption can be easily checked by computing the initial loss value. If it fails, one can either:
\begin{enumerate}
\item Re-initialize the network
\item Use a different initialization scheme
\item Modify the proof to handle the degenerate case (which requires additional technical work)
\end{enumerate}

\item \textbf{Alternative conditions}: If $\xi_2(\cdot)=0$ (second layer frozen) and the initialization has non-zero mass, then non-degeneracy is automatic without this assumption.

\item \textbf{Connection to full-width case}: In the full-width three-layer case, the analogous assumption is $\mathscr{L}(w_1^0,w_2^0,w_3^0) < \E_Z[\mathcal{L}(Y,\varphi_3(0))]$, which has the same interpretation.
\end{itemize}

\subsection{The Convergence Assumption}

The convergence assumption (Assumption~\ref{assump:convergence-lowrank}) is the most technical requirement. It states that the mean-field dynamics converge in a Wasserstein-like sense, with weights that are properly coupled. This assumption is typically verified by:
\begin{itemize}
\item Showing that the loss decreases along trajectories (gradient flow property)
\item Establishing compactness of the trajectory set
\item Using LaSalle's invariance principle or similar stability arguments
\end{itemize}

The low-rank structure enters through the $\max_k$ over $r$ channels, but the overall structure of the assumption remains similar to the full-width case.

\section{Summary}

The proof strategy follows the same path as the full-width case:
\begin{enumerate}
\item Show convergence of the backpropagated signal using the convergence assumption
\item Use coupling arguments to relate $W^*(t)$ to $\bar{W}$
\item Bound differences using the bi-Lipschitz property of $H_2$ (with low-rank factor $rK$)
\item Establish orthogonality using the dense span property
\item Prove non-degeneracy of the limit point
\item Conclude global optimality using convexity
\item Connect the limit using the convergence assumption
\end{enumerate}

The key technical contribution is showing that the low-rank structure (factor $r$) does not break the convergence argument, as long as the universal approximation property is maintained (which is ensured by having $r$ channels with full support).

\section{Extension to $L$ Layers}

We now extend the mean-field framework to $L$-layer low-rank random feature networks. The structure generalizes naturally from the three-layer case.

\subsection{Architecture for $L$ Layers}

Consider an $L$-layer network with the following structure:
\begin{itemize}
\item \textbf{Layer 1}: Random features $\varphi_1(\ip{f_1(C_1)}{X})$ where $f_1(C_1)$ are frozen random vectors and $C_1$ indexes the first-layer channels. \textbf{This layer is not trained.}
\item \textbf{Layers 2 to $L-1$}: Each layer $\ell\in\{2,\dots,L-1\}$ employs a low-rank mixing matrix $L^{(\ell)}$ of rank $r$, producing pre-activations
\[
H_\ell(t,c_\ell;X,W)=\sum_{k=1}^r L^{(\ell)}_{c_\ell,k}\,m_k^{(\ell)}(t;X,W),
\]
where $m_k^{(\ell)}(t;X,W)$ are the $r$ channel-wise partial functions at layer $\ell$.
\item \textbf{Layer $L$ (output)}: Trainable weights $w_L(t,C_L)$ producing the final output.
\end{itemize}

\subsection{Mean-Field ODE System for $L$ Layers}

The mean-field ODE system consists of $(L-2)\times r + 1$ trainable equations:
\begin{itemize}
\item For each intermediate layer $\ell\in\{2,\dots,L-1\}$: $r$ equations for $w_\ell(t,c_\ell,k)$, $k=1,\dots,r$
\item For the output layer $L$: $1$ equation for $w_L(t,c_L)$
\end{itemize}

\textbf{Total:} $(L-2)\times r + 1 \approx rL$ equations for large $L$ (since layer 1 uses frozen random features and is not trained).

For example:
\begin{itemize}
\item $L=3$: $r+1$ equations (as in the three-layer case)
\item $L=4$: $2r+1$ equations
\item $L=5$: $3r+1$ equations
\item General $L$: $(L-2)r+1$ equations
\end{itemize}

For layer $\ell\in\{2,\dots,L-1\}$ and channel $k\in\{1,\dots,r\}$:
\begin{align}
\partial_t w_\ell(t,c_\ell,k) &= -\xi_\ell(t)\,\E_{Z}\!\left[d_L\big(Z;W(t)\big)\,\varphi_\ell'\!\Big(H_\ell(t,c_\ell;X,W(t))\Big)\,B_k^{(\ell)}(t;X)\right], \label{eq:ode-wl}
\end{align}
where $B_k^{(\ell)}(t;X)$ is the channel-$k$ backpropagated signal at layer $\ell$:
\[
B_k^{(\ell)}(t;X)=\E_{C_{\ell+1}}\!\left[L^{(\ell+1)}_{C_{\ell+1},k}\,\varphi_{\ell+1}'\!\Big(H_{\ell+1}(t,C_{\ell+1};X,W(t))\Big)\,w_{\ell+1}(t,C_{\ell+1})\right].
\]

For the output layer $L$:
\begin{align}
\partial_t w_L(t,c_L) &= -\xi_L(t)\,\E_{Z}\!\left[d_L\big(Z;W(t)\big)\,\varphi_L\!\Big(H_L(t,c_L;X,W(t))\Big)\right], \label{eq:ode-wL}
\end{align}
where $H_L(t,c_L;X,W)$ is defined recursively from the previous layers.

\subsection{Recursive Structure of Pre-Activations}

For $\ell\in\{2,\dots,L\}$, the pre-activations satisfy the recursive relation:
\[
H_\ell(t,c_\ell;X,W) = \sum_{k=1}^r L^{(\ell)}_{c_\ell,k}\,m_k^{(\ell)}(t;X,W),
\]
where
\[
m_k^{(\ell)}(t;X,W) = \E_{C_{\ell-1}}\!\left[w_{\ell-1}(t,C_{\ell-1},k)\,\varphi_{\ell-1}\!\Big(H_{\ell-1}(t,C_{\ell-1};X,W(t))\Big)\right].
\]

The model output is:
\[
\hat y(X;W(t)) = \E_{C_L}\big[w_L(t,C_L)\,\varphi_L\big(H_L(t,C_L;X,W(t))\big)\big].
\]

\subsection{Assumptions for $L$ Layers}

\begin{assumption}[Initialization for $L$ layers]\label{assump:init-Llayers}
For each layer $\ell\in\{2,\dots,L\}$, the initial functions satisfy:
\begin{align*}
\sup_{m\ge 1}\frac{1}{\sqrt{m}}\max_{1\le k\le r}\E_{C_\ell}\!\left[|w_\ell^0(C_\ell,k)|^m\right]^{1/m} &\le K \quad \text{for } \ell\in\{2,\dots,L-1\},\\
\sup_{m\ge 1}\frac{1}{\sqrt{m}}\E_{C_L}\!\left[|w_L^0(C_L)|^m\right]^{1/m} &\le K.
\end{align*}
\end{assumption}

\begin{assumption}[Regularity for $L$ layers]\label{assump:regularity-Llayers}
There exists $K\ge 1$ such that:
\begin{itemize}
\item For each $\ell\in\{2,\dots,L\}$, $\xi_\ell$ are bounded on compact intervals and piecewise continuous.
\item For each $\ell\in\{2,\dots,L\}$, $\varphi_\ell,\varphi_\ell'$ are $K$-bounded and $K$-Lipschitz; moreover $\varphi_\ell'$ is non-vanishing for $\ell\in\{2,\dots,L-1\}$.
\item For each $\ell\in\{2,\dots,L-1\}$, the mixing kernels $L^{(\ell)}$ satisfy $\sup_{c_\ell,k}|L^{(\ell)}_{c_\ell,k}|\le K$.
\item All other regularity conditions from Assumption~\ref{assump:regularity-lowrank} hold.
\end{itemize}
\end{assumption}

\begin{assumption}[Convergence for $L$ layers]\label{assump:convergence-Llayers}
There exist functions $\bar{w}_\ell$ for $\ell\in\{2,\dots,L\}$ such that as $t\to\infty$, there exists a coupling $\pi_t$ such that for each intermediate layer $\ell\in\{2,\dots,L-1\}$:
\[
\int (1+|\bar{w}_L|)\prod_{j=\ell+1}^{L-1}(1+|\bar{w}_j|)\max_{1\le k\le r}|\bar{w}_\ell(c_\ell,k)|\,|w_\ell^*(t,c_\ell',k)-\bar{w}_\ell(c_\ell,k)|\,d\pi_t \to 0,
\]
and for the output layer:
\[
\int (1+|\bar{w}_L|)\,|w_L^*(t,c_L')-\bar{w}_L(c_L)|\,d\pi_t \to 0.
\]
\end{assumption}

\subsection{Key Technical Extensions}

\begin{lemma}[Bi-Lipschitz property for $L$ layers]\label{lem:bilipschitz-Llayers}
For any $W'$ and $W''$ satisfying the regularity assumptions, and for any layer $\ell\in\{2,\dots,L\}$:
\[
\big|H_\ell(t,c_\ell;X,W')-H_\ell(t,c_\ell;X,W'')\big| \le K^L\,\|L^{(\ell)}\|_{\infty,1}\,\max_{2\le j\le \ell}\max_{1\le k\le r}\E_{C_j}\big[|w_j'(t,C_j,k)-w_j''(t,C_j,k)|\big],
\]
where the constant $K^L$ comes from the product of Lipschitz constants across layers.
\end{lemma}

\begin{remark}[Scaling with depth]
The key observation is that the low-rank structure at each layer introduces a factor of $r$ (through $\|L^{(\ell)}\|_{\infty,1}\le rK$), but these factors multiply across layers. However, the well-posedness argument remains valid as long as the product $r^L$ is controlled, or more precisely, as long as the $\psi_2$ bounds remain finite. In practice, for fixed $r$ and moderate $L$, this is not an issue.
\end{remark}

\subsection{Proof Strategy for $L$ Layers}

The proof strategy extends naturally from the three-layer case:

\begin{enumerate}
\item \textbf{Well-posedness}: The Picard iteration argument extends to $L$ layers, with the key modification that differences propagate through $L-1$ layers. The contraction constant becomes $C^L$ where $C$ is the single-layer constant, but this is still manageable for fixed $L$.

\item \textbf{Global convergence}: The same argument applies, but now the backpropagated signal $B_k^{(\ell)}$ depends on all subsequent layers. The key bound becomes:
\[
|B_k^{(\ell)*}(t)-B_k^{(\ell)}(\bar{W})| \le K^L(1+|\bar{w}_L|)\prod_{j=\ell+1}^{L-1}(1+|\bar{w}_j|)\Big(|w_L^*(t)-\bar{w}_L| + \sum_{j=\ell+1}^{L-1}\max_k|w_j^*(t)-\bar{w}_j|\Big).
\]

\item \textbf{Convergence assumption}: The assumption must ensure that all layer-wise differences converge, with appropriate weights that account for the depth. The product structure $(1+|\bar{w}_L|)\prod_{j=\ell+1}^{L-1}(1+|\bar{w}_j|)$ ensures that the bound remains finite.
\end{enumerate}

\begin{remark}[Computational complexity]
The number of mean-field equations scales as $(L-2)\times r + 1 \approx rL$ for large $L$, which is linear in both the depth $L$ and the rank $r$. This is much more efficient than the full-width case, which would require $O(N^L)$ equations.
\end{remark}

\end{document}
